\documentclass[12pt]{article}
\usepackage[utf8]{inputenc}
\usepackage[T1]{fontenc}
\usepackage{hyperref}
\title{Suabo Language Resources for K-12 Curriculum Development}
\author{Compiled by Claude Sonnet 5.5 (claude-sonnet-5-5) for Andr\'as Kornai}
\date{}
\begin{document}
\maketitle

\section{Introduction}
Suabo (ISO 639-3 szp), also called Inanwatan, Suabau, Bira, Iagu or Mirabo, is a Papuan language of Southwest Papua, Indonesia. Its speakers call it \emph{nid\'aibo} (``our language''). Estimates are 800 speakers (2004) and about 1,000 in 15 villages along MacCluer Bay (Lourens de Vries, 2010). It is centered on Inanwatan village in South Sorong Regency, and is also spoken in Seget village and in Sorong city. Speakers are mostly older, and the newest generation does not know it. It is usually placed in the South Bird's Head family, though it may form an independent family with Duriankere. A 2024 article by de Vries revisits its position.

\section{Primary Educational Resources}
None found.

\section{Linguistic and Reference Materials}
\begin{itemize}
\item Wikipedia article on the Suabo language: \url{https://en.wikipedia.org/wiki/Suabo_language}
\item Omniglot page on Suabo (\emph{nid\'aibo}): \url{https://omniglot.com/writing/suabo.htm}
\item Leiden University thesis or publication on Inanwatan, found in the results but not opened, so its exact title is unknown: \url{https://scholarlypublications.universiteitleiden.nl/access/item%3A3463909/view}
\item Numerals entry for Suabo (Max Planck Institute for Evolutionary Anthropology): \url{https://lingweb.eva.mpg.de/channumerals/Suabo.htm}
\item A PDF at the Languages of Melanesia site, found in the results but not opened: \url{https://langlxmelanesia.com/_files/ugd/4bb486_4206f5c31f584db68cea182777c1217a.pdf}
\end{itemize}

\section{Religious and Community Resources}
None found.

\section{Digital and Audio Resources}
None found.

\section{Curriculum Development Context}
The language is endangered, and the newest generation does not know it. It is not known to be taught in schools. The documentation I found is linguistic and appears to center on the work of Lourens de Vries and colleagues. A curriculum would need documentation and community involvement first, and no community organizations were found in these searches.
\end{document}